% ----- CHAUPAI 10 -----

\newpage

\subsection*{\deva{दशम चौपाई} (Tenth Chaupai)}
\addcontentsline{toc}{subsection}{\deva{दशम चौपाई} (Tenth Chaupai) — \deva{भीम रूप...}}

\begin{minipage}[t]{0.55\textwidth}
\vspace{0pt}

{\large \linespread{1.5}\selectfont
\deva{भीम रूप धरि असुर सँहारे।}\\
\deva{रामचंद्र के काज सँवारे॥}
\par}

\vspace{0.5em}

{\large \linespread{1.5}\selectfont
\telu{భీమ రూప ధరి అసుర సఁహారే।}\\
\telu{రామచంద్ర కే కాజ సఁవారే॥}
\par}

\vspace{0.5em}

{\large \linespread{1.3}\selectfont
\textit{bhīma rūpa dhari asura sãhāre |}\\
\textit{rāmacandra ke kāja sãvāre ||}
\par}
\end{minipage}%
\hfill
\begin{minipage}[t]{0.42\textwidth}
\vspace{0pt}
\begin{tcolorbox}[anchorbox]
\textbf{Standing Up to Bullies:} Hanuman was inherently peaceful, but when his boundaries were repeatedly broken by toxic forces (\textit{Asura}), he assumed a terrifying, massive form (\textit{Bhima Rupa}) to protect his mission. True strength means being gentle by default, but having the courage and ferocity to stand up decisively to bullies and injustice when necessary.
\vspace{0.5em}\newline When backed into a corner by toxic forces, he unleashed an unapologetic, fierce defense to protect his mission.\newline\null\hfill\href{https://www.valmiki.iitk.ac.in/sloka?field_kanda_tid=5&language=dv&field_sarga_value=47}{\textit{Sundara Kanda, 47:34-37}}
\end{tcolorbox}
\vspace{0.5em}

\end{minipage}

\vspace{0.5em}
\subsubsection*{\textbf{Visual Rhythm Map:}}
\begin{table}[h!]
\centering
\renewcommand{\arraystretch}{1.2}
\noindent\makebox[\textwidth][l]{%
  {%
  \setlength{\tabcolsep}{0pt}%
  \begin{tabularx}{1.0625\textwidth}{|*{17}{>{\centering\arraybackslash\hsize=1.0\hsize}X|}}
  \hline
  \multicolumn{3}{|>{\centering\arraybackslash\hsize=3.0\hsize}X|}{\textbf{\guru{भी}\laghu{म}} \par (3)} &
  \multicolumn{3}{>{\centering\arraybackslash\hsize=3.0\hsize}X|}{\textbf{\guru{रू}\laghu{प}} \par (3)} &
  \multicolumn{2}{>{\centering\arraybackslash\hsize=2.0\hsize}X|}{\textbf{\laghu{ध}\laghu{रि}} \par (2)} &
  \multicolumn{3}{>{\centering\arraybackslash\hsize=3.0\hsize}X|}{\textbf{\laghu{अ}\laghu{सु}\laghu{र}} \par (3)} &
  \multicolumn{6}{>{\centering\arraybackslash\hsize=5.0\hsize}X|}{\textbf{\laghu{सँ}\guru{हा}\guru{रे}} \par (5)} \\
  \hline
  \noalign{\vspace{6pt}}
  \hline
  \multicolumn{6}{|>{\centering\arraybackslash\hsize=6.0\hsize}X|}{\textbf{\guru{रा}\laghu{म}\guru{चं}\laghu{द्र}} \par (6)} &
  \multicolumn{2}{>{\centering\arraybackslash\hsize=2.0\hsize}X|}{\textbf{\guru{के}} \par (2)} &
  \multicolumn{3}{>{\centering\arraybackslash\hsize=3.0\hsize}X|}{\textbf{\guru{का}\laghu{ज}} \par (3)} &
  \multicolumn{6}{>{\centering\arraybackslash\hsize=5.0\hsize}X|}{\textbf{\laghu{सँ}\guru{वा}\guru{रे}} \par (5)} \\
  \hline
  \end{tabularx}%
  }%
}
\vspace{-1.5em}
\end{table}

\vspace{1em}
\textbf{ \deva{सरल-संस्कृतम्}:}\\
 \deva{भवान् भयंकरं (भीमं) रूपं धृत्वा असुरान् संहारितवान्।}\\
 \deva{श्रीरामचन्द्रस्य सर्वाणि कार्याणि भवान् साधितवान् (सफलानि कृतवान्)।}\\


Assuming a massive/terrifying form, you destroyed the demons.\\
And thus, you successfully accomplished the missions of Lord Ramachandra.


\vspace{1em}
\newpage
\textbf{\deva{प्रतिपदार्थः} (Meaning of each word):}
\begin{table}[h!]
\centering
\renewcommand{\arraystretch}{1.2}
\begin{tabularx}{\textwidth}{@{} l l X X @{}}
\toprule
\textbf{Awadhi} & \textbf{Sanskrit} & \textbf{English} & \textbf{Grammar \& Notes} \\
\midrule
\deva{भीम रूप} / \telu{భీమ రూప} / \textit{bhīma rūpa}  &  \deva{भीम-रूपम्}  &  Massive / fearsome form  &   \\
\deva{धरि} / \telu{ధరి} / \textit{dhari}  &  \deva{धृत्वा}  &  Having assumed  &   \\
\deva{असुर} / \telu{అసుర} / \textit{asura}  &  \deva{असुरान्}  &  Demons  &   \\
\deva{सँहारे} / \telu{సఁహారే} / \textit{sãhāre}  &  \deva{संहारितवान्}  &  Destroyed  &   \\
\deva{रामचंद्र के} / \telu{రామచంద్ర కే} / \textit{rāmacandra ke}  &  \deva{श्रीरामचन्द्रस्य}  &  Ramachandra's  & Genitive case marker `ke` \\
\deva{काज} / \telu{కాజ} / \textit{kāja}  &  \deva{कार्याणि}  &  Works/Missions  & Prakrit shift () \\ 
\deva{सँवारे} / \telu{సఁవారే} / \textit{sãvāre}  &  \deva{साधितवान्}  &  Accomplished  &   \\
\bottomrule
\end{tabularx}
\end{table}

\begin{tcolorbox}[poetrybox]
\textbf{The 17-Beat Anomaly \& Bhava-Pradhana:}\\
Mathematically, both lines of this verse sum up to 17 beats due to the heavy weights of words like \deva{संहारे} and \deva{संवारे}. 
But Tulsi das, avoids that problem by using  \deva{सँहारे} and  \deva{सँवारे} ! 
The Hanuman Chalisa is a \textit{\deva{भाव प्रधान}} (emotion foremost) poem.
Singers naturally absorb this 17th beat by slightly eliding or compressing syllables where needed. 
These metrical inconsistencies actually serve as beautiful mnemonic devices that add to the beauty of the singing experience!\\
\textbf{Rhyming Finale:}\\
Notice the brilliant rhyme: by killing the demons (\deva{सँहारे}), he makes Rama's mission a success (\deva{सँवारे}).
\end{tcolorbox}

\newpage
